\chapter{Summary}

\section{Aim of the thesis}

The aim of this thesis was to present the frontend development of WaccApp, a web-based communication platform built on the WhatsApp Business API. The thesis focused on how the technical possibilities of the WhatsApp Business Platform can be transformed into a usable client-side interface. The main objective was not to analyse the backend in detail, but to show how the frontend supports message sending, template-based communication, questionnaire handling, media management, scheduled sending, filtering, and conversation-based data display. 

The project was examined from a frontend-oriented perspective. This means that the thesis focused on the structure of the user interface, the role of the individual HTML views, client-side validation, asynchronous API communication, system state feedback, responsive behaviour, and the practical usability of the implemented functions. The backend was discussed only where necessary to understand how the client-side interface communicates with the server and how the user-facing processes are completed. 

\section{Main frontend results}

The development resulted in a modular, multi-view frontend structure. Instead of placing all functions on a single overloaded page, WaccApp separates the main tasks into dedicated views. The \path{index.html} page serves as the central starting point for sending messages, templates, questionnaires, media files, and scheduled operations. The \path{messages.html}, \path{sent-messages.html}, and \path{conv.html} pages support the display and interpretation of communication history. The \path{questionnaire.html}, \path{template.html}, \path{template-menu.html}, and \path{template-single.html} pages provide the frontend basis for text-based and button-based questionnaire processes, while the \path{filter.html} page supports searching and basic interpretation of stored communication data. 

One of the most important results of the project is that the frontend not only provides access to functions, but also organises them into understandable user workflows. The system supports client-side validation, visible loading, success and error states, and feedback mechanisms that help the user understand what is happening during each operation. This is especially important in a system that communicates with an external platform through API calls, because the user cannot directly see the technical background of message sending, media uploading, questionnaire launching, or scheduled execution. 

The implemented frontend also supports the later interpretation of communication data. Sent and received messages can be displayed in list-based views, while the conversation view presents the interaction with a selected contact in a more natural, chat-like form. The filtering module adds another layer of usability by allowing the user to search communication events and questionnaire-related answers according to different conditions. In this way, WaccApp is not only a sending interface, but also a frontend environment for following and interpreting communication processes. 

\section{Conclusions}

The thesis confirms that the frontend of a WhatsApp Business API-based communication platform must do more than display data or forward form inputs. In WaccApp, the frontend became an active part of the system because it helps the user prepare valid operations, understand system states, avoid common input errors, and follow the results of previous communication. The project, therefore, demonstrates that frontend development has a direct impact on the reliability and practical value of such a communication system. 

The main goals of the thesis were achieved. WaccApp provides a working frontend foundation for message sending, template handling, questionnaire editing and launching, media-related communication, scheduled sending, conversation display, and filtered data retrieval. Although the current version still has several possible development directions, such as a more advanced dashboard, stronger role-based interface logic, improved campaign management, and a more formal component structure, the implemented solution already represents a usable and extensible frontend platform. 

Overall, WaccApp shows that a modular, responsive, and API-driven frontend can make a complex communication system more transparent and easier to use. The value of the project lies not only in the number of implemented functions, but in the way these functions are organised into a coherent user-facing system. The thesis, therefore, concludes that WaccApp is a suitable foundation for further frontend development in the field of WhatsApp Business API-based communication. 
